\section{Simultaneous deterministic choice}\label{sec:accounts}

A single pairing must preserve several estimates at once. We select it
by minimizing a weighted sum of nonnegative statistics, called accounts.
Each account receives its own discount, computed from its exact
conditional expectation. This separates the common choice from the
different growth estimates used to analyze its consequences.
The underlying method of conditional expectations and pessimistic
estimators is classical \cite{Raghavan1988}; see also
\cite[Section~3]{WX2008}. The lemma below
includes the account-specific discount update in the finite averaging
argument.

\begin{lemma}[Discounted accounts]\label{lem:discounted-accounts}
At a state $x$, let $\Omega(x)$ be a nonempty finite set of candidate
transitions with a specified probability distribution.  For each index
$i$ in a finite set let $Z_i(x)\ge1$ be a nonnegative account and let $D_i>0$ be its current
discount.  Fix nonnegative weights $a_i$.  Set
\[
 r_i=\max\left\{1,\frac{\mathbb E[Z_i(x')\mid x]}{Z_i(x)}\right\},
 \qquad D_i'=D_i/r_i.
\]
There is a candidate $x'$ such that
\[
 \sum_i a_iD_i'Z_i(x')\le\sum_i a_iD_iZ_i(x).
\]
In particular, selecting a minimizer of the left side has this property.
If along the selected history
$\mathbb E[Z_i(x')\mid x]\le(1+c_{i,t})Z_i(x)$ with
$c_{i,t}\ge0$, then
\[
 D_{i,t}^{-1}\le\prod_{j<t}(1+c_{i,j})
 \le\exp\left(\sum_{j<t}c_{i,j}\right)
\]
when the initial discount is one.
\end{lemma}

\begin{proof}
For each index,
$D_i'\mathbb EZ_i(x')\le D_iZ_i(x)$ by definition of $r_i$.
Sum with weights $a_i$ and use that a minimum is no larger than its
average.  The discount product follows from $r_i\le1+c_{i,t}$.
\end{proof}

For a normalized nonnegative estimator $Y$ that can vanish, apply the
lemma to $1+Y$. This convention keeps the expectation ratio defined
and preserves every upper bound obtained by extracting $Y$.

An account attached to a surviving label has value $1+\phi_i\ge1$ while
the label is alive and zero after its death.  The same proof applies up
to death; after death its contribution remains zero and its discount
need not be changed.  When a family contains $F$ labels, each is assigned
the fixed weight $1/F$ from the beginning of the phase.  Hence, if the
discounted total is at most $K$ and a surviving label has
$D_i^{-1}\le B$, then $1+\phi_i\le FKB$.  For a normalized $p$th-power
account this costs the factor $(FKB)^{1/p}$, which is explicitly allowed
when the fixed even order $p$ is chosen.
Scalar trace accounts and aggregate low-square accounts have unit
weight; averaging over the $O(\log n)$ spectral parameters would make no
difference.  Polynomial-size fixed denominators are used for the
individual entry and row families, whose extraction uses high moments.

\begin{lemma}[Scalar upper-excess account]\label{lem:scalar-excess}
Let an alive scalar quantity have current value $X$ and upper threshold
$T$.  On its ordinary next transitions suppose the threshold is unchanged
and, for a fixed scale $b>0$, its increment $\Delta$ satisfies
\[
 |\Delta|\le b,\qquad \mathbb E\Delta\le\mu,
 \qquad\mathbb E\Delta^2\le\nu,
 \qquad\mu,\nu\ge0.
\]
For each fixed even integer $p\ge2$ there is $C_p$ such that the account
$Z=1+((X-T)_+/b)^p$ obeys
\[
 \mathbb EZ'\le
 \left[1+C_p\left(\frac\mu b+\frac\nu{b^2}\right)\right]Z.
\]
The same upper bound holds if the ordinary estimates are conditional on
an event, provided that on each other surviving transition the threshold
is changed so that the account does not increase.  A death transition
may set the account to zero.  For two-sided control apply the assertion
to $X$ and $-X$ separately.
\end{lemma}

\begin{proof}
For real $x$ and $|y|\le1$, Taylor's formula on the positive half-line,
including the possible crossing of zero, gives
\[
 (x+y)_+^p\le x_+^p+p x_+^{p-1}y
       +C_p(1+x_+^{p-2})y^2.
\]
Substitute $x=(X-T)/b$ and $y=\Delta/b$, take expectations, and use
$x_+^j\le1+x_+^p$ for $0\le j\le p$.  On other surviving transitions the
account is bounded by its old value, and on death it is zero.  Averaging
these alternatives gives the same bound without conditioning.
\end{proof}

The later pairing phases use this estimate at a prescribed scale.

\begin{corollary}[Positive-excess account at a prescribed scale]
\label{lem:positive-excess-account}
Let $B_s>0$ and $\lambda>0$. Suppose an alive scalar quantity has an ordinary increment $\Delta$
with $|\mathbb E\Delta|\le\mu$,
$\mathbb E\Delta^2\le\nu$, and $|\Delta|\le\lambda B_s$.
For each fixed even $p\ge2$, the upper-excess account for either
sign, measured in units $\lambda B_s$, has expectation factor at most
\begin{equation}
 1+C_p\left(\frac{\mu}{\lambda B_s}
             +\frac{\nu}{\lambda^2B_s^2}\right).
 \label{eq:original-34}
\end{equation}
The threshold, exceptional-transition, and death conventions are
those of Lemma~\ref{lem:scalar-excess}.
\end{corollary}
\begin{proof}
Apply Lemma~\ref{lem:scalar-excess} with $b=\lambda B_s$ separately
to the quantity and its negative.
\end{proof}

We shall take $\lambda=n^{-\zeta}$. If the sum of the drift and variance
terms in \eqref{eq:original-34} is bounded and the total family size is
polynomial, choosing $p$ sufficiently large makes the excess $o(B_s)$.
More generally a permitted loss $n^C$ in the discount is absorbed by the
same extraction, with $p$ chosen after the fixed exponent $C$ has been
determined.

\subsection{The induction and its accumulated error bounds}
Each transition estimate is proved assuming the current displayed caps.
These caps first imply that every allowed candidate precision remains
positive and that all candidate statistics are defined.  The minimizing
choice then gives the next caps by the preceding lemmas and their stated
extraction margins.  This is an induction on the number of transitions: first define all
candidates from the current estimates, then select a candidate, and
finally extract its estimates from the common sum. The largest discount
loss of one account is never charged to another account.

A pathwise drift budget is used only when it can be bounded for every
prefix of the selected history having the asserted caps.  In cleanup,
earlier actual paths are bounded by actual paths in the graph at the end
of that prefix.  The structural accounts control that graph by a
separate discount estimate that does not use the pathwise drift budget.
Thus the induction first obtains the new structural bounds and then uses
them to bound the accumulated discounts of the remaining accounts.

At a phase boundary the accounts may be reinitialized.  Normalization by
a slightly larger small power than the preceding extraction loss makes
their initial total bounded or $n^{o(1)}$.  The number of spectral
parameters is logarithmic, and all label families are polynomial.  An
additive budget $O(\log n)$ is harmless after a normalization by any fixed
positive power; it is not charged as an unscaled multiplicative error.
